
Supposons maintenant @@M0@@ noethérien et @@M1@@ régulier, ce qui entraîne que @@M2@@ est surjectif (3.7). On a vu dans la démonstration de 3.7 que pour tout @@M3@@ il existe @@M4@@ tel que :
@@M5@@
La même propriété reste évidemment vraie après tensorisation par @@M6@@. Par suite, si @@M7@@, il existe @@M8@@ tels que :
@@M9@@
En effet, il suffit de prendre @@M10@@.

On filtre l'anneau @@M11@@ défini en 4.7 en posant :
@@M12@@
Il résulte immédiatement de (4.7.1) que c'est une filtration d'anneau. Pour montrer que pour tout @@M13@@ l'homomorphisme :
@@M14@@
est surjectif, il suffit donc de montrer que si @@M15@@, avec @@M16@@, il existe @@M17@@ tel que @@M18@@.

Or d'après ce qui précède, il existe @@M19@@ tel que @@M20@@ et @@M21@@. D'après 4.7, on a alors :
@@M22@@
Or :
@@M23@@
(conséquence immédiate de (4.7.2)). Comme @@M24@@ est d'augmentation nulle, @@M25@@, et compte tenu de V 5.7,
@@M26@@
d'où le résultat.

